%
% Integranden der Form Polynom/theta
%
\subsection{Integranden der Form $\text{Polynom}/\vartheta$}
Der Summand $\varphi$ in
\eqref{buch:intalg:wurzel:eqn:summe}
lässt sich auf ein Integral nur von $1/\vartheta$ reduzieren,
wie der folgende Satz zeigt.

\begin{satz}
\label{buch:intalg:wurzel:satz:polynomtheta}
Sei $\varphi\in\mathbb{Q}[x]$ ein Polynom mit rationalen Koeffizienten
und $\vartheta=\!\sqrt{ax^2+bx+c}$.
Dann gibt es ein Polynom $\psi\in\mathbb{Q}[x]$ mit 
\begin{equation}
\int\frac{\varphi}{\vartheta}\,dx
=
\psi\vartheta
+
A\int\frac{1}{\vartheta}\,dx,
\label{buch:intalg:wurzel:eqn:polytheta}
\end{equation}
wobei $A$ eine Konstante ist.
\end{satz}

\begin{proof}
Wenn es das Polynom $\psi$ und die Konstante $A$ gibt, dann 
muss auch die Ableitung der Gleichung~\eqref{buch:intalg:wurzel:eqn:polytheta} 
erfüllt sein, also
\begin{equation}
\frac{\varphi}{\vartheta}
=
\psi'\vartheta + \psi\vartheta'+A\frac{1}{\vartheta}.
\label{buch:intalg:wurzel:eqn:polytheta'}
\end{equation}
Die Ableitung von $\vartheta$ ist
\[
\vartheta'
=
\frac{2ax+b}{2\vartheta}
=
\frac{ax+b/2}{\vartheta}.
\]
Eingesetzt in 
\eqref{buch:intalg:wurzel:eqn:polytheta'} und mit $\vartheta$ multipliziert
erhält man die Gleichung
\begin{align*}
\varphi
&=
\psi'\vartheta^2 + \psi (ax+b/2) + A
\\
%\varphi
&=
\psi' (ax^2+bx+c) + \psi (ax+b/2) + A.
\end{align*}
Um diese Differentialgleichung für die Funktion $\psi$ zu lösen, verwenden
wir Ansätze der Form
\[
\renewcommand{\arraycolsep}{1pt}
\begin{array}{rcrcrcrcrcccrcr}
\varphi&=& \varphi_mx^m &+& \varphi_{m-1}x^{m-1} &+&   \varphi_{m-2}x^{m-2} &+&   \varphi_{m-3}x^{m-3} &+& \ldots &+& \varphi_1 x &+& \varphi_0\phantom{.}\\
\psi   &=&              & &    \psi_{m-1}x^{m-1} &+&      \psi_{m-2}x^{m-2} &+&      \psi_{m-3}x^{m-3} &+& \ldots &+&    \psi_1 x &+&    \psi_0\phantom{.}\\
\psi'  &=&              & &                      & & (m-1)\psi_{m-1}x^{m-2} &+& (m-2)\psi_{m-2}x^{m-3} &+& \ldots &+&   2\psi_2 x &+&    \psi_1.
\end{array}
\]
Durch Ausmultiplizieren und Koeffizientenvergleich entstehen die Gleichungen
\[
\renewcommand{\arraycolsep}{1pt}
\def\i#1{\rlap{$\scriptstyle #1$}\phantom{m-2}}
\begin{array}{lcrcrcrcrcrcrcr}
\varphi_{m  } &=& (m-1)a\psi_{m-1}   & &                    & &                    &+& a\psi_{m-1}   & &                   & &   \\
\varphi_{m-1} &=& (m-2)a\psi_{m-2}   &+& (m-1)b\psi_{m-1}   & &                    &+& a\psi_{m-2}   &+& (b/2)\psi_{m-1} & &   \\
\varphi_{m-2} &=& (m-3)a\psi_{m-3}   &+& (m-2)b\psi_{m-2}   &+& (m-1)c\psi_{m-1}   &+& a\psi_{m-3}   &+& (b/2)\psi_{m-2}   & &   \\
\varphi_{m-3} &=& (m-4)a\psi_{m-4}   &+& (m-3)b\psi_{m-3}   &+& (m-2)c\psi_{m-2}   &+& a\psi_{m-4}   &+& (b/2)\psi_{m-3}   & &   \\
              &\vdots &              & &                    & &                    & &               & &                   & &   \\
\varphi_{3}   &=&     2a\psi_{\i{2}} &+&     3b\psi_{\i{3}} &+&     4c\psi_{\i{4}} &+& a\psi_{\i{2}} &+& (b/2)\psi_{\i{3}} & &   \\
\varphi_{2}   &=&      a\psi_{\i{1}} &+&     2b\psi_{\i{2}} &+&     3c\psi_{\i{3}} &+& a\psi_{\i{1}} &+& (b/2)\psi_{\i{2}} & &   \\
\varphi_{1}   &=&                    & &      b\psi_{\i{1}} &+&     2c\psi_{\i{2}} &+& a\psi_{\i{0}} &+& (b/2)\psi_{\i{1}} & &   \\
\varphi_{0}   &=&                    & &                    & &      c\psi_{\i{1}} & &               &+& (b/2)\psi_{\i{0}} &+& A \\
\end{array}
\]
oder zusammengefasst
\[
\renewcommand{\arraycolsep}{1pt}
\def\i#1{\rlap{$\scriptstyle #1$}\phantom{m-2}}
\begin{array}{lcrcrcrcrcr}
\varphi_{m  } &=& \phantom{(}m\phantom{\mathstrut-0)}a\psi_{m-1}
                                     & &                             & &                    & &   \\
\varphi_{m-1} &=& (m-1)a\psi_{m-2}   &+&  ((2(m-1)+1)b/2)\psi_{m-1}  & &                    & &   \\
\varphi_{m-2} &=& (m-2)a\psi_{m-3}   &+&  ((2(m-2)+1)b/2)\psi_{m-2}  &+& (m-1)c\psi_{m-1}   & &   \\
\varphi_{m-3} &=& (m-3)a\psi_{m-4}   &+&  ((2(m-3)+1)b/2)\psi_{m-3}  &+& (m-2)c\psi_{m-2}   & &   \\
              &\vdots &              & &                             & &                    & &   \\
\varphi_{3}   &=&     3a\psi_{\i{2}} &+& 7b/2\phantom{)}\psi_{\i{3}} &+&     4c\psi_{\i{4}} & &   \\
\varphi_{2}   &=&     2a\psi_{\i{1}} &+& 5b/2\phantom{)}\psi_{\i{2}} &+&     3c\psi_{\i{3}} & &   \\
\varphi_{1}   &=&      a\psi_{\i{0}} &+& 3b/2\phantom{)}\psi_{\i{1}} &+&     2c\psi_{\i{2}} & &   \\
\varphi_{0}   &=&                    & &  b/2\phantom{)}\psi_{\i{0}} &+&      c\psi_{\i{1}} &+& A.\\
\end{array}
\]
Daraus kann man Rekursionsformeln für die Koeffizienten $\psi_k$ ableiten:
\begin{align*}
\psi_{m-1}
&=
\frac{1}{ma}\varphi_m
\\
\psi_{m-2}
&=
\frac{1}{(m-1)a}\biggl(
\varphi_{m-1}
-
\frac{(2(m-1)+1)b}{2}\psi_{m-1}
\biggr)
\\
\psi_{m-3}
&=
\frac{1}{(m-2)a}\biggl(
\varphi_{m-2}
-
\frac{(2(m-1)+1)b}{2}\psi_{m-2}
-
(m-1)c\psi_{m-1}
\biggr)
\\[-5pt]
&\phantom{a}\vdots
\\[-3pt]
\psi_{k}
&=
\frac{1}{(k+1)a}
\biggl(
\varphi_{k+1}
-
\frac{(2(k+1)+1)b}{2}\psi_{k+1}
-
(k+2)
\psi_{k+2}
\biggr)
\\[-5pt]
&\phantom{a}\vdots
\\[-3pt]
\psi_0
&=
\frac{1}{a}\biggl(
\varphi_1
-
\frac{3b}{2}\psi_1
-
2c\psi_2
\biggr)
\\
A
&=
\varphi_0 -(b/2)\psi_0 - c\psi_1.
\end{align*}
Damit sind der Reihe nach die Koeffizienten
$\psi_{m-1},\psi_{m-2},\dots,\psi_2,\psi_1,\psi_0,A$
eindeutig festgelegt, was die Existenz des Polynoms $\psi$ und der
Konstanten $A$ beweist.
\end{proof}

Der Satz besagt, dass der erste Summand in~\eqref{buch:intalg:wurzel:eqn:summe}
zu einem Integranden führt, dessen Integral sich auf das Integral von
$1/\vartheta$ reduzieren lässt.
Bei der Berechnung des zweiten Summanden treten noch weitere Arten
solcher spezieller Integrale auf.
Diese werden in Abschnitt~\ref{buch:intalg:wurzel:pol:subsection:ktheta}
auf ein Integral mit einem Integranden der Form $1/\vartheta$ zurückgeführt
und im Abschnitt~\ref{buch:intalg:wurzel:subsection:stammfunktiontheta}
berechnet.
